% Apêndice — detalhamento, responsabilidade a responsabilidade, da
% correspondência entre as responsabilidades do harness e a literatura.
% Histórico: antiga subseção 2.1.3 da Fundamentação; migrado para a Metodologia
% em 17/10/2026 (decisão D5d); em 07/11/2026 virou subsubseção da atividade 3;
% em 11/11/2026 passou a apêndice, por decisão do autor na revisão da
% Metodologia (análise de literatura já feita não é atividade planejada). O
% conteúdo dos parágrafos não foi alterado; só a abertura foi adaptada.
% Resumo no Quadro qua:harness-liu (Cap. 3). Bases: A28 e
% context/planning/mapeamento-harness-liu2025.md.
\apendice{Delimitação das responsabilidades do \textit{harness}}
\label{ap:delimitacao}

Este apêndice examina, responsabilidade a responsabilidade, o que o catálogo de
\citeonline{liu2025} cobre e o que fica a cargo de outras fontes ou de decisão do
projeto, segundo o resumo do Quadro~\ref{qua:harness-liu}. Atende ao primeiro
objetivo específico, que delimita as responsabilidades da camada com base em
padrões arquiteturais publicados, e é o ponto de partida do projeto da camada
(Seção~\ref{subsec:at-projeto}).

\paragraph{Responsabilidade 1.} O \textit{Agent adapter} converte as mensagens do
agente para o formato exigido por ferramentas externas e vice-versa, e o
\textit{Prompt/response optimiser} padroniza \textit{prompts} e respostas por
meio de modelos de texto (\textit{templates}) \cite{liu2025}. Nenhum dos dois
trata da troca do provedor do modelo: o primeiro lida com as interfaces das
ferramentas, e o segundo, com a estrutura dos \textit{prompts}. É, porém, a
responsabilidade cujo efeito o experimento mede. Para ancorá-la, recorre-se à
noção de \textit{middleware}, a camada reutilizável que absorve a heterogeneidade
entre plataformas \cite{issarny2007}, discutida na
Seção~\ref{sec:abstracao}. A correspondência é de propósito, e não de projeto:
o conceito explica por que uma camada intermediária absorve a heterogeneidade,
mas não indica como projetá-la para provedores de modelos, e essa transposição é
do autor.

\paragraph{Responsabilidade 2.} O \textit{Agent adapter} invoca e gerencia as
interfaces das ferramentas e devolve ao modelo os resultados da execução
\cite{liu2025}, o que cobre parte da função. O \textit{Tool/agent registry} trata
da seleção entre muitas ferramentas e agentes, problema que pouco se coloca
quando o conjunto de ferramentas é pequeno e conhecido de antemão, como se prevê
neste trabalho. Já os padrões \textit{One-shot model querying} e
\textit{Incremental model querying} descrevem quantas vezes o modelo é consultado
para gerar um plano; relacioná-los ao ciclo de ferramentas é uma analogia do
autor.

\paragraph{Responsabilidade 3.} O \textit{Passive goal creator} recupera da
memória o histórico de conversas, os artefatos em edição e as ferramentas usadas
recentemente, e os anexa ao \textit{prompt} do usuário; o \textit{Incremental
model querying} divide o contexto em vários \textit{prompts} para contornar o
limite da janela de contexto; e o \textit{Retrieval augmented generation}
enriquece a consulta com conhecimento externo \cite{liu2025}. Nenhum deles trata
da gestão do estado da interação como responsabilidade de uma camada. Para o
estado, recorre-se a \citeonline{chang2025}, cujo SagaLLM integra o padrão
transacional \textit{Saga} a memória persistente e compensação automática, para
tratar a perda de contexto e a ausência de reversão no planejamento multiagente.
O escopo multiagente está fora deste trabalho, e o que se retém é a separação
entre dois problemas: preservar contexto e estado ao longo das etapas, e
desfazer ou compensar etapas já executadas. A segunda, por ser compensação
semântica, não é tratada pela camada; o tratamento de falhas da
responsabilidade 4 limita-se às de transporte e de protocolo.\pendencia{Confirmar
com o orientador que a recuperação de falhas se limita a transporte e protocolo
e que a compensação semântica fica explicitamente fora do escopo.}

\paragraph{Responsabilidade 4.} A tolerância a falhas (\textit{fault tolerance})
aparece no catálogo apenas como força e benefício da cooperação baseada em
papéis, no sentido de substituir agentes indisponíveis, e a seção 5.3 cita a
identificação e o tratamento de exceções em tempo real como trabalho em
andamento dos autores \cite{liu2025}. A literatura de arquitetura de software
trata o mesmo problema por outro ângulo, o de lidar com serviços remotos sujeitos
a falha: \textit{Retry} e \textit{Circuit Breaker} são padrões de resiliência
usados com essa finalidade em microsserviços \cite{aderaldo2025}, já analisados
por meio de modelos \cite{mendonca2020} e avaliados experimentalmente
\cite{sedghpour2022,aderaldo2025}. Do ponto de vista da aplicação, o provedor de
IA é um serviço remoto, sujeito às mesmas classes de falha transitória, como
tempo esgotado, erro de conexão e limite de uso. Transpor esses padrões das
chamadas entre microsserviços para as chamadas a provedores de modelos é uma
analogia do autor: as fontes não tratam de IA. Há ainda uma indicação útil para a
questão deste trabalho: os próprios padrões de resiliência introduzem atrasos de
invocação, e o efeito desses atrasos sobre o desempenho da aplicação ainda era
pouco estudado \cite{aderaldo2025}.

\paragraph{Responsabilidade 5.} Também não há padrão entre os 18. Ao discutir a
integração do catálogo com padrões existentes, \citeonline{liu2025} citam padrões
de IA responsável publicados em outra obra do mesmo grupo, entre eles o
\textit{black box}, que coleta as entradas e saídas do agente em tempo de
execução para servir de evidência em processos de auditoria. Esses padrões não
integram o catálogo e não são tratados neste trabalho. Para o que se entende
aqui por registro estruturado de eventos, adota-se a documentação do
OpenTelemetry: um registro (\textit{log}) é um registro textual com marca de
tempo, estruturado (formato recomendado pela documentação) ou não, com metadados
opcionais, e o seu modelo de dados prevê campos de identificação do rastreio
(\textit{trace}) e do segmento (\textit{span}) da requisição (\textit{TraceId} e
\textit{SpanId}), que associam o registro à requisição em cujo processamento ele
foi emitido \cite{opentelemetrylogs}. Trata-se de documento técnico, não revisado
por pares, usado apenas para fixar o conceito de registro estruturado e
correlacionado.

\paragraph{Risco da própria camada.} O catálogo registra, por fim, um risco que
interessa diretamente a este trabalho. Entre as desvantagens do
\textit{Tool/agent registry}, os autores apontam que ele pode se tornar uma
solução de aprisionamento a fornecedor (\textit{vendor lock-in}) e um ponto único
de falha \cite{liu2025}. O contexto é outro, o do registro de ferramentas, e não
o da integração com provedores de modelos, mas o argumento vale para a camada
estudada aqui: abstrair o provedor para reduzir a dependência pode criar uma
dependência nova, na própria camada.
